\documentclass[12pt]{article}
\usepackage{notes-project}

\newcommand{\abbottexercise}[1]{%
  \emph{Source:} Stephen Abbott, \emph{Understanding Analysis}, Second
  Edition, Exercise #1.\par\smallskip
}

\begin{document}

\begin{center}
  {\LARGE\bfseries Tutorial 1}\\[0.25em]
  {\large AMA3707 Real Analysis}\\[0.15em]
\end{center}

\section{Sets and Functions}

\begin{exc}{abbott-1-2-3}
\abbottexercise{1.2.3}
\end{exc}

\textbf{(a)} If $A_1 \supseteq A_2 \supseteq A_3 \supseteq \cdots$ and
every $A_n$ is infinite, then $\bigcap_{n=1}^{\infty} A_n$ is infinite.

Given that $A_1 \subseteq A_2 \subseteq A_3 \subseteq \cdots$ are all containing infinite number of elements. 

This is a nested sequence that 
\[
\bigcap_{n=1}^{m} A_n = A_m
\]

Notice that for all sets $A_n$, it contains infinite number of elements. Meaning that $m \to \infty$, $A_m$ will still contain infinite number of elements. 

Also, 
\[
\bigcap_{n=1}^{\infty} A_n = \lim_{m \to \infty}{\bigcap_{n=1}^{m} A_n} = A_m
\] 

So, 
\[
\left| \bigcap_{n=1}^{\infty} A_n \right| = \left| A_m \right| = \infty 
\] 

Statement (a) is true. 

\medskip
\noindent\textbf{(b)} If $A_1 \supseteq A_2 \supseteq A_3 \supseteq \cdots$
and every $A_n$ is finite and nonempty, then
$\bigcap_{n=1}^{\infty} A_n$ is finite and nonempty.

By the same argument, 
\[
\bigcap_{n=1}^{\infty} A_n = \lim_{m \to \infty}{\bigcap_{n=1}^{m} A_n} = A_m, \quad 0 < |A_m| < \infty
\] 

So, 
\[
\bigcap_{n=1}^{\infty} A_n = A_m \quad \text{is finite and nonempty.}
\] 

Statement (b) is true. 

\medskip
\noindent\textbf{(c)} $A \cap (B \cup C) = (A \cap B) \cup C$.

By \xref[membership under classical logic]
{lec01::rem:membership-classical-logic}, every element \(x\) in the chosen
universe satisfies exactly one of the following eight conditions:
\begin{enumerate}[nosep] 
\item $x \in A$ and $x \in B$ and $x \in C$
\item $x \in A$ and $x \in B$ and $x \notin C$
\item $x \in A$ and $x \notin B$ and $x \in C$
\item $x \in A$ and $x \notin B$ and $x \notin C$
\item $x \notin A$ and $x \in B$ and $x \in C$
\item $x \notin A$ and $x \in B$ and $x \notin C$
\item $x \notin A$ and $x \notin B$ and $x \in C$
\item $x \notin A$ and $x \notin B$ and $x \notin C$
\end{enumerate} 

Define 
\[D= A \cap (B \cup C), \quad E = (A \cap B) \cup C\] 

Denote the above 8 conditions as $P_1, P_2, \ldots, P_8$. 

By using definition of \xref{lec01::def:intersection} and \xref{lec01::def:union}, we can see that 
\begin{itemize}[nosep]
\item $\forall x \in P_1 \implies x \in D$ and $x \in E$
\item $\forall x \in P_2 \implies x \in D$ and $x \in E$
\item $\forall x \in P_3 \implies x \in D$ and $x \in E$
\item $\forall x \in P_4 \implies x \notin D$ and $x \notin E$
\item $\forall x \in P_5 \implies x \notin D$ and $x \in E$
\item $\forall x \in P_6 \implies x \notin D$ and $x \notin E$
\item $\forall x \in P_7 \implies x \notin D$ and $x \in E$
\item $\forall x \in P_8 \implies x \notin D$ and $x \notin E$
\end{itemize}

So, $D = P_1 \cup P_2 \cup P_3$ and $E = P_1 \cup P_2 \cup P_3 \cup P_5 \cup P_7$. 

$\exists x, x \in E \land x \notin D$ such as $x \in P_5$ or $x \in P_7$, meaning that $D \subset E$. 

According to the definition of \xref[set equality]{lec01::def:set-equality}, $D \neq E$. 

Statement (c) is false. 

\medskip
\noindent\textbf{(d)} $A \cap (B \cap C) = (A \cap B) \cap C$.



\medskip
\noindent\textbf{(e)}
$A \cap (B \cup C) = (A \cap B) \cup (A \cap C)$.

\begin{exc}{abbott-1-2-7}
\abbottexercise{1.2.7}
\end{exc}

% Write your answer below this line.

\begin{exc}{abbott-1-2-9}
\abbottexercise{1.2.9}
\end{exc}

% Write your answer below this line.

\section{Suprema and Infima}

\begin{exc}{abbott-1-3-3}
\abbottexercise{1.3.3}
\end{exc}

% Write your answer below this line.

\begin{exc}{abbott-1-3-4}
\abbottexercise{1.3.4}
\end{exc}

% Write your answer below this line.

\begin{exc}{abbott-1-3-11}
\abbottexercise{1.3.11}
\end{exc}

% Write your answer below this line.

\end{document}
